\begin{center}
    {\large\bfseries 4. Не меньше; не больше; хотя бы один}
\end{center}

\footnotesize

\begin{enumerate}
    \item Иллюзия спора

    Прочитайте четыре диалога. Во всех случаях второй собеседник уверен, что опровергает первого. Однако в трёх случаях из четырёх он допускает логическую ошибку (подменяет тему, путает свойства или даже случайно подтверждает слова собеседника).    
    Найдите единственный диалог с опровержением.

    \vspace{2mm}

    \textit{Олимпиада.} Артём: «В нашей команде каждый участник решил хотя бы одну задачу».\\
    Боря: «Ты не прав. Седьмую задачу не решил вообще никто.»

    \vspace{2mm}

    \textit{Копилка.} Вера: «Среди любых четырёх монет из этой копилки найдётся не менее двух фальшивых».\\
    Глеб: «Нет. Мне не повезло, и среди четырёх монет, которые я вытащил, оказалось целых три фальшивых.»

    \vspace{2mm}

    \textit{Спортивная секция.} Даша: «В нашей секции каждый мальчик подтягивается не больше двух раз».\\
    Егор: «Это неправда. В нашей секции есть три мальчика, которые не могут подтянуться ни разу!»

    \vspace{2mm}

    \textit{Шарики в коробке.} Женя: «В этой коробке шариков любого цвета не более четырёх».\\
    Стёпа: «Ты ошиблась: вот, я достал оттуда пять синих шариков».


    \item Равносильные утверждения.

    К каждому утверждению найдите равносильное.

    1. Не все задачи олимпиады имеют одинаковую сложность.
    \begin{enumerate}
        \item Все задачи олимпиады имеют разную сложность.
        \item В олимпиаде найдутся хотя бы две задачи разной сложности.
        \item В олимпиаде есть хотя бы одна очень простая и хотя бы одна очень сложная задача.
        \item Большинство задач олимпиады имеют различную сложность.
    \end{enumerate}

    \vspace{0.5em}
    2. Среди любых пяти яблок в ящике найдётся хотя бы одно червивое.
    \begin{enumerate}
        \item В ящике червивых яблок больше, чем хороших.
        \item В ящике не может быть больше четырёх хороших яблок.
        \item В ящике лежит ровно четыре хороших яблока.
        \item В ящике есть хотя бы пять червивых яблок.
    \end{enumerate}
    \vspace{0.8em}


    \item Расположите утверждения так, чтобы получилась цепочка следствий:

    А. В кармане есть не менее 5 конфет.

    Б. В кармане есть хотя бы 1 конфета.
    
    В. В кармане ровно 20 конфет.

    Г. В кармане не менее 15 конфет.


\item Баба Яга в своей избушке на курьих ножках завела сказочных животных. Все они, кроме двух, — Говорящие Коты; все, кроме двух, — Мудрые Совы; остальные — Усатые Тараканы. Сколько обитателей в избушке у Бабы Яги?

\item Итак, мама воскликнула — «Чудеса!», и сразу же мама, папа и дети отправились в зоомагазин. «Но здесь больше пятидесяти снегирей, как мы выберем», — чуть не заплакал младший брат, увидев снегирей. «Не волнуйся», — сказал старший, — «их меньше пятидесяти». «Главное,» — сказала мама, — «что здесь есть хотя бы один!» «Да, забавно,» — подытожил папа, — «из трех ваших фраз только одна соответствует действительности». Сможете ли Вы сказать, сколько снегирей было в магазине, зная, что снегиря мне купили?
\end{enumerate}
